%% kapa: exam=1 round=26 year=2015 period=1 form=A subject=civil_law subject_ko=민법 seq=21 number=21
%% source: Q-Net 감정평가사 기출문제 (cst00309) artlSeq=5138001 file=extracted/1008/5138001/제26회_감정평가사 1차_기출/제26회_감정평가사 1차_기출/1교시_감평_A형.pdf
%% conversion: kordoc 4.24.51 -> tools/gen1.py (auto); review=proofread-sonnet; flags=none
\begin{question}{21}
부동산 등기에 관한 설명으로 옳은 것은? (다툼이 있으면 판례에 따름)
\begin{choices}
  \item 등기에 공신력이 인정된다.
  \item 지상권 설정등기가 불법 말소된 경우 그 지상권은 소멸한다.
  \item 동일인 명의로 보존등기가 중복된 경우 후등기가 무효이다.
  \item 멸실된 건물의 보존등기를 그 대지 위에 신축한 건물의 보존등기로 유용할 수 있다.
  \item 매매를 원인으로 하여 甲에서 乙 앞으로 마쳐진 소유권이전등기에 대해 甲이 매매의 부존재를 이유로 그 말소를 청구하는 경우, 乙은 등기의 추정력을 주장할 수 없다.
\end{choices}
\end{question}
\answer{3}
